\documentclass[11pt]{article}

% Math
\usepackage{amsmath, amssymb, amsfonts, mathtools}

% Page layout
\usepackage[margin=1in]{geometry}

% Optional useful packages
\usepackage{physics}      % \bra, \ket, \dv, \pdv, etc.
\usepackage{bm}           % bold math symbols
\usepackage{bbm}
\usepackage{enumitem}     % better lists
\usepackage{graphicx}     % Required for inserting images
\usepackage{subcaption}   % Required for the subfigure environment
\usepackage{hyperref}     % links
\hypersetup{
    colorlinks=true,       % false: boxed links; true: colored links
    linkcolor=blue,        % color of internal links (sections, pages, etc.)
    citecolor=green,       % color of citation links to bibliography
    % filecolor=magenta,     % color of file links
    urlcolor=blue,         % color of external website links
    % allcolors=blue       % uncomment this if you want ALL links to use one color
}

\title{Weekly Update}
\author{Reza Asakereh}
\date{Last Update: \today}

\begin{document}

\maketitle

\section{Main Findings}
\subsection{DB-Sorter}
\textbf{quick background.} DB-sorter uses the sorting characteristic of DB flow to find ground state or excited states of a hamiltonian $H$ using a diagonal matrix $D$ that has an engineered order in its diagonal elements (usually increasing). 
\\

\textbf{Findings:}
\begin{itemize}
    \item I chatted with ChatGPT and it annihilated my work so far. Apparently what I have needs more work. Here is a list of suggestions:
    \begin{itemize}
        \item Only proposing a few "good" $D$'s is not that good. Let's see more theoretical work on how a (near-)optimum $D$ look like
        \item There is not enough work done regarding excited state, despite the claim is it can also address them
        \item More work on sorting convergence can be done besides simply fixed point
        \item Also compare against methods like subspace-VQE and folded spectrum
    \end{itemize}
    \item Optimum $D$ can be roughly found using Lagrange multipliers and a variational circuit. However my method so far is exponentially expensive. See \hyperref[sec:optimum-D]{3.1.1}
\end{itemize}


\section{Plan For Next Week(s)}
\subsection{DB-sorter}
\begin{itemize}
    \item Finding a more efficient way of finding optimum $D$.
    \item More analysis on excited states and sorting behaviour.
    \item Comparison with subspace-VQE and folded spectrum.
\end{itemize}


\section{Details}
\subsection{DB-sorter}
\subsubsection{Best D}
\label{sec:optimum-D}
Optimally, we want the path from $H(0)$ to $H(\infty)$ be the shortest, which means minimizing 
$$
\int^\infty_0 \lVert dH \rVert
$$
However, on the quantum circuit, we only implement the first few steps. Therefore in a greedy way, we want to maximize the change from $H_0$ to $H_1$ or in other words:
\begin{align}
    \max \lVert dH \rVert &= \max_D tr([[D, H], H]^2) \\
    &=\max_D tr(6H^2DH^2D - 8H^3DHD + 2H^4D^2) \\
    &= \max_{d_1, \dots, d_n} \sum_{k=1}^n{d_k^2(6H^2[k,k]H^2[k,k] - 8H^3[k,k]H[k,k] + 2H^4[k,k])} \\
    &= \max_{d_1, \dots, d_n} \sum_{k=1}^n{d_k^2m_k}
\end{align}

Our implementation of bracket relies on $e^{-s[D, H]} = e^{\sqrt{s}D}e^{\sqrt{s}H}e^{-\sqrt{s}D}e^{-\sqrt{s}H}$. In this implementation we have two observations: a) adding a constant value to all elements of $D$ would not affect the final results and b) multiplying a constant to all elements of $D$ can be equivalent of changing our step size. These two observations let us assume that all elements of $D$ are in $[0, 1]$. We try two scenarios:

Case I: $d_1 = 0 \le d_2 \le \dots \le d_n = 1$ and Case II: $d_1 = 0 \le d_2, d_3, \dots, d_{n-1} \le d_n = 1$.
Both of these can be solved using KKT Lagrange multipliers.

For case I, the problem description would be:
\begin{align}
    &\max \sum_{k=1}^n{d_k^2m_k} \\
    \text{s.t } d_k &\ge d_{k-1}: k\ge 2\\
    d_1 &= 0, d_n = 1
\end{align}
which results in the Lagrangian:
\begin{align}
    L(\{d_k\}^n_1, \{\lambda_k\}^n_2) = \sum_{k=1}^n{d_k^2m_k} - \sum_{k=2}^n{\lambda_k(d_{k-1} - d_k)}
\end{align}
solution:
\begin{align}
    2d_km_k - \lambda_{k+1} + \lambda_k &= 0 \Rightarrow d_k = \frac{\lambda_{k+1} - \lambda_k}{2m_k}\\
    \lambda_k(d_k - d_{k-1}) &= 0 \\
    \lambda_k &\ge 0
\end{align}

Considering the complementary slackness, we can test different number of starting $k$'s for which $\lambda_k = 0$. If the first $r$ Lagrange multipliers are zero, from above equations we will have 
$$
d_1, \dots, d_{r-2} = 0, d_{r-1} = \frac{\lambda_r}{2m_r}
$$
With induction and complementary slackness, you can show that (except for an edge case), all feasible solutions have $\lambda_{k>r} = \alpha \lambda_r, \alpha > 0$, therefore $d_{k \ge r-1} = 1$. This means that the solution to this optimization is a diagonal matrix with the first $r$ elements equal to 0 and the rest equal to 1.


For case II, the problem description would be:
\begin{align}
    &\max \sum_{k=1}^n{d_k^2m_k} \\
    \text{s.t } 0 &\le d_k \le 1\\
    d_1 &= 0, d_n = 1
\end{align}
which results in the Lagrangian:
\begin{align}
    L(\{d_k\}^n_1, \{\lambda_k\}^n_1, \{\gamma_k\}^n_1) = \sum_{k=1}^n{d_k^2m_k} - \sum_{k=1}^n{\lambda_k(d_k - 1) + \sum_{k=1}^n{\gamma_kd_k}}
\end{align}
solution:
\begin{align}
    2d_km_k - \lambda_k + \gamma_k &= 0 \Rightarrow d_k = \frac{\lambda_k - \gamma_k}{2m_k}\\
    \lambda_k(d_k - 1) &= 0 \\
    \gamma_kd_k &= 0 \\
    \lambda_k, \gamma_k &\ge 0
\end{align}
Using complementary slackness, we should have all elements arbitrarily $0$ or $1$.

Both cases boil down to exponentially many candidates to check the maximum in. Also computing the trace in equation 2 requires exponentially many Pauli terms.


\end{document}
